%%%%%%%%%%%%%%%%%%%%%%%%%%%%%%%%%%%%%%%%%%%%%%%%%%%%%%%%%%%%%%%%%%%%%%
% Lớp 8 | Chương IV. Định lí Thalès | Đường trung bình của tam giác
% Dạng 8H4?16-2: Chứng minh trung điểm, song song nhờ đường trung bình
%%%%%%%%%%%%%%%%%%%%%%%%%%%%%%%%%%%%%%%%%%%%%%%%%%%%%%%%%%%%%%%%%%%%%%


% nguồn: dan_tu_AI (nhập 16/11/2026)
\begin{bt}%[8H4H16-2]%[ma:6f39268e]
	Cho tam giác $ABC$. Gọi $M$ là trung điểm của $AB$. Qua $M$ kẻ đường thẳng song song với $BC$ cắt $AC$ tại $N$. Chứng minh $N$ là trung điểm của $AC$ và $MN=\dfrac12BC$.
	\loigiai{
		Vì $MN\parallel BC$ nên theo định lí Thalès: $\dfrac{AN}{NC}=\dfrac{AM}{MB}=1$ (vì $AM=MB$).\\
		Suy ra $AN=NC$, tức $N$ là trung điểm của $AC$.\\
		Khi đó $MN$ là đường trung bình của tam giác $ABC$, nên $MN=\dfrac12BC$.
	}
\end{bt}

% nguồn: dan_tu_AI (nhập 16/11/2026)
\begin{bt}%[8H4H16-2]%[ma:2edc09f9]
	Cho tứ giác $ABCD$. Gọi $E$, $F$, $G$, $H$ lần lượt là trung điểm của $AB$, $BC$, $CD$, $DA$. Chứng minh tứ giác $EFGH$ là hình bình hành.
	\loigiai{
		$EF$ là đường trung bình của tam giác $ABC$ nên $EF\parallel AC$ và $EF=\dfrac12AC$.\\
		$HG$ là đường trung bình của tam giác $ADC$ nên $HG\parallel AC$ và $HG=\dfrac12AC$.\\
		Suy ra $EF\parallel HG$ và $EF=HG$.\\
		Vậy tứ giác $EFGH$ là hình bình hành.
	}
\end{bt}

% nguồn: dan_tu_AI (nhập 16/11/2026)
\begin{bt}%[8H4H16-2]%[ma:f466fc45]
	Cho tam giác $ABC$. Gọi $D$, $E$ lần lượt là trung điểm của $AB$, $AC$.\\
	a) Chứng minh tứ giác $BDEC$ là hình thang.\\
	b) Nếu tam giác $ABC$ cân tại $A$ thì tứ giác $BDEC$ là hình gì? Vì sao?
	\loigiai{
		a) $DE$ là đường trung bình của tam giác $ABC$ nên $DE\parallel BC$.\\
		Vậy tứ giác $BDEC$ là hình thang.\\
		b) Tam giác $ABC$ cân tại $A$ nên $\widehat{B}=\widehat{C}$.\\
		Hình thang $BDEC$ ($DE\parallel BC$) có hai góc kề đáy $BC$ bằng nhau, nên $BDEC$ là hình thang cân.
	}
\end{bt}

% nguồn: dan_tu_AI (nhập 16/11/2026)
\begin{bt}%[8H4H16-2]%[ma:7c45f467]
	Cho tam giác $ABC$ vuông tại $A$. Gọi $M$, $N$, $P$ lần lượt là trung điểm của $AB$, $AC$, $BC$. Chứng minh tứ giác $AMPN$ là hình chữ nhật.
	\loigiai{
		$PN$ là đường trung bình của tam giác $ABC$ nên $PN\parallel AB$, tức $PN\parallel AM$.\\
		$PM$ là đường trung bình của tam giác $ABC$ nên $PM\parallel AC$, tức $PM\parallel AN$.\\
		Tứ giác $AMPN$ có các cặp cạnh đối song song nên là hình bình hành.\\
		Hình bình hành $AMPN$ có $\widehat{MAN}=90^\circ$ nên là hình chữ nhật.
	}
\end{bt}

% nguồn: dan_tu_AI (nhập 16/11/2026)
\begin{bt}%[8H4H16-2]%[ma:a52b8d1e]
	Cho tam giác $ABC$ có hai đường trung tuyến $BD$ và $CE$ cắt nhau tại $G$. Gọi $I$, $K$ lần lượt là trung điểm của $GB$, $GC$. Chứng minh tứ giác $DEIK$ là hình bình hành.
	\loigiai{
		$D$, $E$ là trung điểm $AC$, $AB$ nên $ED$ là đường trung bình của tam giác $ABC$: $ED\parallel BC$, $ED=\dfrac12BC$.\\
		$I$, $K$ là trung điểm $GB$, $GC$ nên $IK$ là đường trung bình của tam giác $GBC$: $IK\parallel BC$, $IK=\dfrac12BC$.\\
		Suy ra $ED\parallel IK$ và $ED=IK$.\\
		Vậy tứ giác $DEIK$ là hình bình hành.
	}
\end{bt}

% nguồn: dan_tu_AI (nhập 16/11/2026)
\begin{bt}%[8H4V16-2]%[ma:7e47fe47]
	Cho hình thang $ABCD$ ($AB\parallel CD$) có $AB=6\text{ cm}$, $CD=10\text{ cm}$. Gọi $M$ là trung điểm của $AD$. Qua $M$ kẻ đường thẳng song song với $AB$ cắt $BD$ tại $I$ và cắt $BC$ tại $N$.\\
	a) Chứng minh $I$ là trung điểm của $BD$ và $N$ là trung điểm của $BC$.\\
	b) Tính độ dài $MN$.
	\loigiai{
		a) Tam giác $DAB$ có $M$ là trung điểm của $AD$ và $MI\parallel AB$ nên $I$ là trung điểm của $BD$.\\
		Tam giác $BDC$ có $I$ là trung điểm của $BD$ và $IN\parallel DC$ (vì $MN\parallel AB\parallel CD$) nên $N$ là trung điểm của $BC$.\\
		b) $MI$ là đường trung bình của tam giác $DAB$ nên $MI=\dfrac{AB}{2}=3\text{ cm}$.\\
		$IN$ là đường trung bình của tam giác $BDC$ nên $IN=\dfrac{CD}{2}=5\text{ cm}$.\\
		Vì $I$ nằm giữa $M$ và $N$ nên $MN=MI+IN=3+5=8\text{ (cm)}$.
	}
\end{bt}
